\clearpage\thispagestyle{empty}
\pdfbookmark[0]{Sources for New Elements of Pallas}{pallaselements1810oct-sources}
\begin{center}\textbf{Sources for New Elements of Pallas}\end{center}
\noindent This AI translation follows \textit{Neue Pallaselemente} in \href{https://gdz.sub.uni-goettingen.de/id/PPN236007467}{\textit{Werke}, vol.~VI} (Göttingen, 1874), p.~320. Item XLV of \href{https://archive.org/download/bub_gb_YCI4AAAAMAAJ/page/n4.jpg}{\textit{Monatliche Correspondenz}, vol.~XXII}, October 1810, names Gauss on p.~\href{https://archive.org/download/bub_gb_YCI4AAAAMAAJ/page/n403.jpg}{400} and prints his ephemeris on pp.~\href{https://archive.org/download/bub_gb_YCI4AAAAMAAJ/page/n404.jpg}{401}, \href{https://archive.org/download/bub_gb_YCI4AAAAMAAJ/page/n405.jpg}{402} and \href{https://archive.org/download/bub_gb_YCI4AAAAMAAJ/page/n406.jpg}{403}. The digitized University of Michigan copy is available from the Internet Archive.
\par\smallskip\noindent The collected version preserves the report and orbital elements, followed by a bracket describing the ephemeris. The first-person plural is the voice of the report, which names Gauss in the third person as the person who calculated the ephemeris. The ascending-node row, the comma-shaped leader in the inclination row and the angular marks follow the collected version. Translation and typesetting produced by AI.
\clearpage
